\documentclass[12pt]{article}
\usepackage[utf8]{inputenc}
\usepackage[T1]{fontenc}
\usepackage{hyperref}
\title{Dubuporo Language Resources for K-12 Curriculum Development}
\author{Compiled by Claude Sonnet 5.5 (claude-sonnet-5-5) for Andr\'as Kornai}
\date{}
\begin{document}
\maketitle

\section{Introduction}
No search result mentions the name ``Dubuporo''. The Glottocode yang1298 and the ISO code ynl point to the language called Yangulam, a Rai Coast (Madang) language of Madang Province, Papua New Guinea, in Nuclear Trans New Guinea. Wikipedia gives about 400 speakers from the 2000 census, and ASJP gives 910. ``Dubuporo'' may be an alternative or local name, but the searches did not confirm this. The language is not known to be taught in schools. Little material was found.

\section{Primary Educational Resources}
None were found in the search results.

\section{Linguistic and Reference Materials}
\begin{itemize}
\item Wikipedia article on Yangulam: \url{https://en.wikipedia.org/wiki/Yangulam_language}
\item Glottolog entry yang1298: \url{https://glottolog.org/resource/languoid/id/yang1298.bigmap.html}
\item ASJP word list, drawn from Z'graggen (1980): \url{https://asjp.clld.org/languages/YANGULAM}
\item Trans New Guinea database entry: \url{https://transnewguinea.org/language/yangulam?page=6}
\end{itemize}

\section{Religious and Community Resources}
None were found in the search results.

\section{Digital and Audio Resources}
A search summary mentioned an open-access PARADISEC collection of Yangulam materials from Don Daniels. No URL for it appeared in the results, so it is not linked here and should be checked at the PARADISEC catalogue.

\section{Curriculum Development Context}
The language appears to be vigorous among a few hundred speakers but has almost no published educational material. Curriculum work would start from the archival recordings and word lists. The community name and orthography should be confirmed with speakers.
\end{document}